\documentclass{jsarticle}
\usepackage{amsmath}
\usepackage{amssymb}
\begin{document}
\title{}
\author{}
\date{}
\maketitle

\begin{flushleft}
    順子お姉さんへ
\end{flushleft}

       ベータ関数へと導くゼータ関数についての場の理論となる
       サーストン・ペレルマン多様体から、小学校時代の姉同士の公園での
       素因数分解の教えあいも、
    聞いていたことと、中学校時代の姉同士の河合塾がいいよの話と、
    高等学校での、私に教えてくれた数学での学び、
    あとは、福山市での本での学びを、結実した代数がトポロジーへと
    つながり、ゼータ関数がシャノンの公式とコルモゴロフ方程式、
    非可換代数多様体へと、全部は、すべて、姉たちの教えも兼ねて導けました。
  

        私に骨を埋めずに、ちゃんとした彼氏を先祖供養をきっかけに見つけて、
    結婚されてください。いい子たちをもらえると思います。
    この子たちに、ピアノの速読を教えて、この速読で数学、物理、化学を
    学んでもらって、成長を期待してください。絶対、いい子たちに育ちます。


        大変な中、この手紙を見てくれて、本当にありがとうございます。



       高齢出産の可能性があり、私の手相も高齢出産の手相になっていて、
    姉も高齢出産と養子縁組の可能性らしく、お互いがんばりましょう。



       あと、姉は、水星は、伝令の神であり、順から、この名前の由来は、
    伝令と侍、仏教から来ている。木星は、多津基くんから、神々の神、
    土星は、知の神、火星は、快楽の神、金星は、美の神、天王星は、
    優しさの神、冥王星は、冥土の神であり、占いは、易学の数秘術から
    創られているのを姉は、知っているし、本当にすみませんでした。


\begin{flushright}
	雅旭より
\end{flushright}
\end{document}
